\documentclass{article}
\title{総合問題集 — 図書館全課程ドリル}
\author{TeX 図書館}

\begin{document}

\section{この問題集のために}

この本は「読む」本ではなく「解く」本です。図書館の教科書(数学・統計・投資の数学・プログラミング 4 分野)を読み終えた人の\textbf{総点検}を目的に、32 問を集めました。

使い方は 3 つのかしこさがあります。

\begin{enumerate}
\item \textbf{通しで解く}: 各分野を順に解いて、弱点の分野を特定する
\item \textbf{教科書に戻る}: 解けなかった問題の「解答の指針」から、対応する教科書の節へ戻って読み直す
\item \textbf{1 週間後に再解く}: 解けなかった問題だけをもう一周する(間隔反復)
\end{enumerate}

すべての問題に「解答の指針」が付いており、答えは折りたたまれています。答えを見る前に、紙かエディタで手を動かすのが前提の設計です。

\section{数学の問題(8 問)}

\textbf{問題 1.}\ \(\displaystyle \lim_{x \to 0} \frac{\sin 3x}{x} \) を求めよ。

\textbf{解答の指針:}\ \(\frac{\sin 3x}{x} = 3 \cdot \frac{\sin 3x}{3x}\) で、\(3x \to 0\) のとき \(\frac{\sin 3x}{3x} \to 1\)。よって値は 3(『大学数学入門』第 6 節の重要極限)。

\textbf{問題 2.}\ \(f(x) = x^3 - 6x^2 + 9x + 1\) の極大値と極小値を求めよ。

\textbf{解答の指針:}\ \(f'(x) = 3(x-1)(x-3)\)。\(x = 1\) で極大値 5、\(x = 3\) で極小値 1(第 6 節)。

\textbf{問題 3.}\ \(\displaystyle \int_0^{\pi/2} \sin x \, dx \) を求めよ。

\textbf{解答の指針:}\ 不定積分は \(-\cos x\)。値は \( [-\cos x]_0^{\pi/2} = 0 - (-1) = 1 \)(第 7 節)。

\textbf{問題 4.}\ \(A = \begin{pmatrix} 2 & 1 \\ 1 & 2 \end{pmatrix}\) の固有値をすべて求めよ。

\textbf{解答の指針:}\ 固有方程式 \((2-\lambda)^2 - 1 = 0\) より \(\lambda = 1, 3\)。対角化も可能(第 5 節)。

\textbf{問題 5.}\ 問題 4 の行列 \(A\) に対して、\(A^n\) の固有値の和を求めよ(\(n\) は自然数)。

\textbf{解答の指針:}\ \(A^n\) の固有値は \(1^n, 3^n\) なので和は \(1 + 3^n\)。トレースとの一致で検算できる(第 5 節)。

\textbf{問題 6.}\ ベクトル \(\begin{pmatrix} 3 \\ 4 \end{pmatrix}\) と同じ向きの単位ベクトルを求めよ。

\textbf{解答の指針:}\ 大きさが 5 なので \(\frac{1}{5}\begin{pmatrix} 3 \\ 4 \end{pmatrix} = \begin{pmatrix} 3/5 \\ 4/5 \end{pmatrix}\)(第 2 節)。

\textbf{問題 7.}\ \(f(x, y) = x^2 + y^2 - 2x + 4y + 5\) の最小値を求めよ。

\textbf{解答の指針:}\ \(\frac{\partial f}{\partial x} = 2x - 2 = 0\), \(\frac{\partial f}{\partial y} = 2y + 4 = 0\) より \((1, -2)\)。\(D = f_{xx}f_{yy} - f_{xy}^2 = 4 > 0\) かつ \(f_{xx} = 2 > 0\) で極小。値は \(f(1, -2) = 0\)(第 8 節)。

\textbf{問題 8.}\ \(|r| < 1\) のとき \(\displaystyle \sum_{k=1}^{\infty} k r^{k-1} = \frac{1}{(1-r)^2}\) を、等比級数の和を \(x\) で微分することで導け。

\textbf{解答の指針:}\ \(\sum_{k=0}^\infty r^k = \frac{1}{1-r}\)。両辺を \(r\) で微分すると \(\sum_{k=1}^\infty k r^{k-1} = \frac{1}{(1-r)^2}\)(第 9 節の総合演習と同型)。

\section{統計の問題(8 問)}

\textbf{問題 9.}\ 平均 50、標準偏差 10 の母集団から 25 個の標本を取る。標本平均が 53 を超える確率を求めよ。

\textbf{解答の指針:}\ 標準誤差は \(10/\sqrt{25} = 2\)。\(z = (53 - 50)/2 = 1.5\) より \(P(Z > 1.5) \approx 0.067\)(『統計学入門』第 7 節)。

\textbf{問題 10.}\ データ \([2, 4, 6, 8, 10]\) の分散(母分散)と標準偏差を求めよ。

\textbf{解答の指針:}\ 平均 6、偏差 \((-4, -2, 0, 2, 4)\)、\(s^2 = \frac{16 + 4 + 0 + 4 + 16}{5} = 8\)、標準偏差 \(2\sqrt{2}\)(第 2 節)。

\textbf{問題 11.}\ 相関係数 \(r = -0.6\) のとき、回帰直線の傾きは正か負か。また決定係数は?

\textbf{解答の指針:}\ 傾き \(b = r \cdot s_y / s_x\) は負。\(R^2 = 0.36\) で 36\% を説明(第 10 節)。

\textbf{問題 12.}\ 100 回のトスで表が 60 回出た。表の出る確率を 0.5 と仮定したときの z 統計量を求めよ(二項分布の標準偏差は \(\sqrt{np(1-p)}\))。

\textbf{解答の指針:}\ \(np = 50\), \(\sqrt{np(1-p)} = \sqrt{25} = 5\)。\(z = (60 - 50)/5 = 2.0\)。両側で p 値約 4.6\% なので 5\% 水準で棄却——コインは疑わしい(第 9 節)。

\textbf{問題 13.}\ 有病率 2\% の病気の検査で感度 90\%、偽陽性率 5\% とする。陽性の人が実際に病気である確率を求めよ。

\textbf{解答の指針:}\ ベイズ: \(\frac{0.9 \times 0.02}{0.9 \times 0.02 + 0.05 \times 0.98} = \frac{0.018}{0.067} \approx 26.9\%\)(第 4 節)。

\textbf{問題 14.}\ \(n = 16\), \(\bar{x} = 40\), \(s = 12\) のとき母平均の 95\% 信頼区間を求めよ(自由度 15 の t 値は 2.131)。

\textbf{解答の指針:}\ \(40 \pm 2.131 \times 12/4 = 40 \pm 6.39\)、よって \([33.6,\ 46.4]\)(第 8 節)。

\textbf{問題 15.}\ 「p 値 0.03 は帰無仮説が正しい確率 3\% を意味する」——この文の誤りを説明せよ。

\textbf{解答の指針:}\ p 値は「\(H_0\) のもとでデータ以上に極端な結果の確率」で条件の向きが逆。\(H_0\) の正しさの確率は事前確率と検出力に依存する(ベイズの定理と第 9 節の注意)。

\textbf{問題 16.}\ ポアソン分布 Po(3) に従う事故件数で「1 日あたり事故が 5 件以上」になる確率の見積もり方を述べよ(計算式まで)。

\textbf{解答の指針:}\ \(P(X \ge 5) = 1 - \sum_{k=0}^{4} e^{-3} \frac{3^k}{k!}\)(第 6 節)。\(e^{-3}(1 + 3 + 4.5 + 4.5 + 3.375) \approx 0.815\) より約 18.5\%。

\section{投資の数学の問題(6 問)}

\textbf{問題 17.}\ 勝率 40\%、ペイアウト比 \(b = 2\) のトレードの期待値係数を求めよ。

\textbf{解答の指針:}\ \(E = 0.4 \times 3 - 1 = 0.2\)。賭け金の 20\% が期待値(『投資の数学』第 3 節)。

\textbf{問題 18.}\ 問題 17 のケリー比を求めよ。

\textbf{解答の指針:}\ \(f^* = E / b = 0.2 / 2 = 0.1\)。資金の 10\% が最適賭け幅(第 6 節)。

\textbf{問題 19.}\ 資金の 20\% を賭け続けて 5 連敗した。資金の残率を求めよ。

\textbf{解答の指針:}\ \(0.8^5 \approx 32.8\%\)(第 7 節)。固定比率方式では順序に依存せずこの値に戻る道がある。

\textbf{問題 20.}\ オッズ 2.4 の試合を勝率 50\% と推定した。賭けるべきか、期待値で判断せよ。

\textbf{解答の指針:}\ \(b = 1.4\)、\(E = 0.5 \times 2.4 - 1 = 0.2 > 0\) で賭け候補。ただし推定の信頼性次第で半ケリー以下に(第 5・6 節)。

\textbf{問題 21.}\ 「負けるたびに賭け金を 2 倍にする」戦略で 7 連敗したときの累積賭け額を求めよ。

\textbf{解答の指針:}\ \(1 + 2 + 4 + \cdots + 64 = 127\)。資金の限界で続けられず、それまでの勝ちが一瞬で消える(第 9 節)。

\textbf{問題 22.}\ 勝率 55\%、\(b = 1\) のゲームのケリー比と、半ケリーの賭け割合を求めよ。

\textbf{解答の指針:}\ \(f^* = 0.1\)、半ケリーは 5\%(第 6 節)。小さい賭け幅でも長期成長はプラスで、破産耐性が大幅に上がる。

\section{プログラミングの問題(10 問)}

\textbf{問題 23.}\ \([x ** 2 for x in range(1, 5)]\) の評価結果を答えよ。

\textbf{解答の指針:}\ \([1, 4, 9, 16]\)(『Python 入門』第 4 節)。

\textbf{問題 24.}\ 辞書 \texttt{\{"a": 1, "b": 2\}} のキーの一覧を得るコードを書け。

\textbf{解答の指針:}\ \texttt{list(d.keys())} またはループで \texttt{for k in d}(第 3 節)。

\textbf{問題 25.}\ ファイルを例外が起きても確実に閉じたい。使う構文と理由を述べよ。

\textbf{解答の指針:}\ \texttt{with open(...) as f:} — ブロック終了時に自動 close、例外時も保証(第 7 節)。

\textbf{問題 26.}\ \texttt{Counter("mississippi").most_common(1)} の評価結果を答えよ。

\textbf{解答の指針:}\ \texttt{[("i", 4)]}(i が 4 回で最多)。\texttt{("s", 4)} と同数だが先に現れた要素が優先される(第 6 節)。

\textbf{問題 27.}\ TypeScript で \texttt{const xs = [1, 2, 3]} のとき \texttt{xs.map(x => x / 2)} の戻り値の型を答えよ。

\textbf{解答の指針:}\ \texttt{number[]}。map は要素型から新しい配列型を作る(『TypeScript 入門』第 7 節)。

\textbf{問題 28.}\ 判別可能ユニオンで「ローディング中 / 成功 / 失敗」の状態を型で表現せよ。

\textbf{解答の指針:}\ 判別子 \texttt{state} をもつ 3 つのケースのユニオンにする — \texttt{isLoading} と \texttt{data?} の組合せより不正状態を排除できる(第 7・10 節)。

\textbf{問題 29.}\ SQL で「ユーザーごとの合計購入額を多い順に」取得するクエリの骨格を書け(orders テーブル)。

\textbf{解答の指針:}\ \texttt{SELECT user\_id, SUM(amount) AS total FROM orders GROUP BY user\_id ORDER BY total DESC;}(『SQL・データベース入門』第 3 節)。

\textbf{問題 30.}\ LEFT JOIN の後で \texttt{WHERE o.id IS NULL} と書くと何が得られるか。

\textbf{解答の指針:}\ 結合相手が存在しない左側の行 = 「注文のないユーザー」など(第 4 節)。

\textbf{問題 31.}\ 10 億件のソート済みデータから 1 件を二分探索する際の最大比較回数を求めよ。

\textbf{解答の指針:}\ \(\lceil \log_2 10^9 \rceil = 30\) 回(『アルゴリズムとデータ構造』第 4 節)。

\textbf{問題 32.}\ 「ある配列の中から合計が target になる組合せがあるか」を \(O(2^n)\) でなく \(O(n \cdot target)\) 程度で解く発想を述べよ。

\textbf{解答の指針:}\ 動的計画法——到達できる合計値の集合を順に更新し、状態の重複をまとめて処理する(第 9 節)。

\section{おわりに — 解けなかったら戻る}

32 問を解き終えたら、次の対応表で弱点の教科書へ戻るのが最短です。

\begin{center}
\begin{tabular}{|l|l|}
\hline
分野 & 戻るべき教科書 \\
\hline
数学(問題 1〜8) & 大学数学入門 \\
統計(問題 9〜16) & 統計学入門 \\
投資(問題 17〜22) & 投資の数学 \\
プログラミング(問題 23〜32) & Python 入門 / TypeScript 入門 / SQL・データベース入門 / アルゴリズム \\
\hline
\end{tabular}
\end{center}

解けなかった問題は、本棚の「学習カード」でも同じテーマのデッキが用意されています(各本のサイドバー → 🎴 学習カード)。問題集で弱点を見つけ、カードで穴を埋め、1 週間後に問題集へ戻る——このサイクルが図書館の学習設計の完成形です。

\end{document}
